\chapter{Variations of Knapsack Cryptosystems}

\section{Weaknesses of the Merkle--Hellman Scheme}

The classical Merkle--Hellman cryptosystem is based on a clear trapdoor idea.
An easy private superincreasing subset-sum instance is transformed into a
public instance that should look hard~\cite{MerkleHellman1978}. However, in the
basic construction, this hiding transformation is not strong enough to remove
all useful structure. As a result, the public key does not behave like a hard
subset-sum instance, and this opens the door to efficient cryptanalysis.

A major result was given by Shamir, who showed a polynomial-time
attack against the basic Merkle--Hellman scheme \cite{Shamir1982}. The attack
demonstrates that, even without knowing the private sequence directly, an
equivalent private structure can be reconstructed from public information. This
breaks the intended trapdoor asymmetry and shows that the original scheme is
insecure.

\subsection{Shamir's Polynomial-Time Attack}

Shamir showed that the basic Merkle--Hellman cryptosystem can be broken in
polynomial time.~\cite[pp.~145--152]{Shamir1982} The attack is not based on exhaustive search
over ciphertexts, but on structure that remains visible in the public key. This
means that the weakness lies in the key construction itself.

Importantly, breaking the scheme does not require reconstructing the exact
private sequence chosen during key generation. It is sufficient to recover an
equivalent trapdoor that allows the public knapsack to be transformed into an
easy superincreasing instance. Once such a trapdoor is available, ciphertexts
can be decrypted efficiently, which is already enough to break the
system.~\cite[pp.~86--87]{OdlyzkoRiseFall}

Shamir's attack is not a separate ciphertext attack performed from
scratch for every message. Instead, it is a key-recovery attack with an
initial polynomial-time computation on the public key. After this
precomputation, the attacker can decrypt later ciphertexts in essentially the
same way as the legitimate recipient. This is what makes the attack
particularly strong against the basic Merkle--Hellman construction.~\cite[pp.~77--78]{OdlyzkoRiseFall}

To see where this structure comes from, let $u \equiv r^{-1} \pmod{m}$. Then
each public weight $b_i$ satisfies
\[
    u b_i \equiv w_{\pi(i)} \pmod{m},
\]
where $\pi$ denotes the possible secret permutation of the private sequence.
Equivalently, for each $i$ there exists an integer $k_i$ such that
\[
u b_i - k_i m = w_{\pi(i)}.
\]
Dividing by $m b_i$ gives
\[
\frac{u}{m} - \frac{k_i}{b_i} = \frac{w_{\pi(i)}}{m b_i}.
\]
For those indices corresponding to the smallest hidden weights, the right-hand
side is very small, so the fractions $k_i/b_i$ are all close to the same
unknown ratio $u/m$. This is the key observation behind Shamir's attack: the
public key still contains enough arithmetic information to approximate the
hidden modular parameters and reconstruct an equivalent trapdoor.~\cite[pp.~84--85]{OdlyzkoRiseFall}

Shamir's main contribution was to show that these approximation relations can
be exploited in polynomial time. From a small collection of such relations, the
attacker can obtain a sufficiently accurate approximation to the hidden ratio
$u/m$ and then construct another pair of modular parameters that transforms the
public knapsack into a superincreasing one. These recovered parameters do not
have to coincide with the original secret ones. It is enough that they define
an equivalent trapdoor that allows efficient decryption. In this way, the
attack turns the structural weakness of the public key into a full
polynomial-time break of the classical
scheme.~\cite[pp.~145--152]{Shamir1982}\cite[pp.~85--86]{OdlyzkoRiseFall}

\subsection{Low-Density Attacks and Lattice Methods}

Shamir's attack showed that the classical Merkle--Hellman scheme can fail
because its hidden trapdoor structure is not concealed well enough. Later
analysis showed that this is only part of the picture. Another important
class of attacks does not try to recover the original trapdoor at all, but
targets the public subset-sum instance directly when it has low density. In
these conditions, the public knapsack may be solvable much more efficiently than by
exhaustive search, which broadens the weakness from a single construction flaw
to a more general structural problem in knapsack
cryptography.~\cite[pp.~82--83]{OdlyzkoRiseFall}\cite{LagariasOdlyzko1985}

A standard measure used in this context is the density of the subset-sum
instance, usually defined as
\[
    d=\frac{n}{\log_2(\max_i a_i)},
\]
where \(a_1,\dots,a_n\) are the public weights. When \(d\) is small, the public
weights are large compared with the number of message bits. Then the subset
sums are spread over a larger range of values. Because of this, the vector that
represents the correct \(0\)-\(1\) solution can stand out as a short vector in
the constructed lattice.

The lattice embedding can be described as follows. Suppose that the ciphertext
is \(C\), where
\[
    C=\sum_{i=1}^{n} x_i a_i
\]
for an unknown message vector \(x=(x_1,\dots,x_n)\in\{0,1\}^n\). Consider the
lattice generated by the rows of the matrix
\[
    B=
    \begin{pmatrix}
        1 & 0 & \cdots & 0 & -a_1\\
        0 & 1 & \cdots & 0 & -a_2\\
        \vdots & \vdots & \ddots & \vdots & \vdots\\
        0 & 0 & \cdots & 1 & -a_n\\
        0 & 0 & \cdots & 0 & C
    \end{pmatrix}.
\]
The negative values \(-a_i\) are used so that the last coordinate measures the
difference between the ciphertext and a selected subset sum. If \(v_i\) denotes
the \(i\)-th row of this matrix and \(x\) is the correct solution, then
\[
    \sum_{i=1}^{n} x_i v_i + v_{n+1}
    =
    (x_1,x_2,\dots,x_n, C-\sum_{i=1}^{n}x_i a_i)
    =
    (x_1,x_2,\dots,x_n,0).
\]
Thus the solution appears inside the lattice as a vector with only \(0\) and
\(1\) in the first coordinates and zero in the last coordinate. This makes it
short compared with many other lattice vectors.

For example, take \(a=(7,11,19)\) and \(C=30\). The solution is
\(x=(0,1,1)\), since \(30=11+19\). The corresponding basis vectors include
\[
    v_2=(0,1,0,-11),\quad v_3=(0,0,1,-19),\quad v_4=(0,0,0,30),
\]
and therefore
\[
    v_2+v_3+v_4=(0,1,1,0).
\]
This example is small and not meant as a secure parameter choice, but it shows
the cancellation idea behind the lattice construction.

The important point is that the solution vector is usually not one of the
original basis vectors. It is hidden as an integer combination of them. Lattice
basis reduction tries to replace the original basis by another basis of the same
lattice, but with shorter and less dependent vectors. The number of basis
vectors does not change; only the vectors themselves are replaced by integer
combinations of the old ones. The LLL algorithm performs this process in
polynomial time. Internally, it uses Gram--Schmidt orthogonalization to measure
the lengths and directions of basis vectors, and then decides when to reduce a
vector by another one or when to swap two basis vectors. It does not guarantee
that the shortest vector will always be found, but in many low-density
subset-sum instances the solution vector is short enough to be revealed by the
reduced basis.~\cite{LenstraLenstraLovasz1982,LagariasOdlyzko1985}

Lagarias and Odlyzko gave a general theoretical framework for this type of
attack. They showed that sufficiently low-density subset-sum instances can
often be solved in polynomial time by lattice-based methods, without any need
to reconstruct the original trapdoor. For knapsack cryptography, this was an
important shift in perspective: even if the private structure is not recovered,
the public instance may still be weak simply because its parameters place it in
a regime where lattice reduction is effective.~\cite{LagariasOdlyzko1985}

Low-density attacks complement Shamir's result. Instead of
recovering an equivalent secret structure, they show that the public knapsack
itself may already be directly solvable.

\subsection{Why NP-Completeness Does Not Guarantee Security}

The NP-completeness of subset-sum is a statement about the worst-case
complexity of the general decision problem. A cryptosystem, however, does not
publish arbitrary instances of this problem. It generates highly structured
instances that must remain easy for the legitimate user to solve with the help
of secret information. Because of this, worst-case hardness alone does not say
enough about the security of the concrete public keys used in practice. As
Diffie and Hellman already observed, complexity theory gives no indication of
the difficulty of a particular knapsack instance.~\cite[p.~654]{diffie1976newDirections}
The history of knapsack cryptography illustrates this limitation clearly. The
Merkle--Hellman scheme was motivated by the NP-completeness of subset-sum, yet
Shamir's structural attack and later low-density attacks showed that the
generated public instances could be much easier than arbitrary subset-sum
instances. The central question is therefore not whether subset-sum is hard in
general, but whether the specific family of instances produced by the
cryptosystem avoids exploitable structure. In the case of knapsack
cryptosystems, this turned out to be very difficult to
achieve.~\cite[pp.~81--83]{OdlyzkoRiseFall}

\section{Selected Variations of Knapsack Cryptosystems}

The attacks discussed above did not immediately end interest in knapsack
cryptography. They motivated attempts to modify or redesign the original
Merkle--Hellman construction. Some variants stayed close to the original idea
and tried to hide the superincreasing sequence more thoroughly, for example by
using permutations or repeated modular transformations. Other proposals changed
the construction more substantially.

\subsection{Modified Merkle--Hellman Schemes}

The first modifications of the Merkle--Hellman idea stayed close to the
original construction. They did not replace the superincreasing trapdoor,
but tried to disguise it more strongly. Merkle and Hellman already noted
that the code breaker's task could be complicated by scrambling the order of
the public weights and by adding different multiples of the modulus to
them.~\cite[p.~526]{MerkleHellman1978} In modular notation, the public
weights still satisfy a relation of the form
\[
    b_i \equiv r w_{\pi(i)}
    \pmod{m}, 
\]
where $W=(w_1,\dots,w_n)$ is the private superincreasing sequence, $m$ is a
secret modulus, $r$ is invertible modulo $m$, and $\pi$ is a possible secret
permutation. These changes make the public sequence less directly readable, but
they do not remove the basic relation between the public and private weights.

The main proposal was to repeat the modular multiplication step several times
instead of applying it only once. In a $Y$-times iterated Merkle--Hellman
knapsack, one starts with an easy private sequence and produces the public
weights through a chain of modular transformations.    More explicitly, if
$a^{(0)}$ denotes the original easy sequence, then the $j$-th stage produces a
new sequence by
\[
    a_i^{(j)} \equiv r_j a_i^{(j-1)} \pmod{m_j}, \qquad j=1,\dots,Y.
\]
The final public sequence is $a^{(Y)}$, possibly after a secret
permutation. For decryption, the legitimate receiver applies the modular
inverses in the opposite order and eventually obtains a subset-sum instance
over the original easy sequence. Thus the trapdoor idea is the same as in the
basic Merkle--Hellman scheme, but repeated several times.

Each stage is intended to
further obscure the structure inherited from the previous one. The hope was
that repeated disguising would hide the original trapdoor more effectively than
the single-iteration construction.~\cite[pp.~342--344]{Brickell1984}


\subsection{Brickell's Attack on Iterated Knapsacks}

Brickell's attack targets the iterated construction itself, not only individual
ciphertexts. The attack uses lattice basis reduction on selected subsets of the
public weights. The reduced lattice basis is expected to contain short vectors
that come from the hidden modular relations between the successive
transformations. Brickell calls the useful short vectors desirable vectors, and
uses them to recover information about the hidden order and structure of the
knapsack.~\cite[pp.~344--348]{Brickell1984}

From these short vectors, the attacker can reconstruct enough information to
obtain a sequence that behaves like a superincreasing sequence for most
positions. The remaining positions can then be handled separately. In this
sense, the attack does not need to recover the exact original private key. It
is sufficient to recover an equivalent structure that makes decryption
efficient.~\cite[pp.~349--355]{Brickell1984}

Although Brickell did not claim a proof that the attack succeeds for every
parameter choice, he reported successful attacks on examples with $n=100$ and
with $Y=5$, $10$, and $20$, which showed that iteration was not only
theoretically questionable but also practically fragile.~\cite[pp.~343,
356--357]{Brickell1984} At the same time, Brickell noted that each iteration
reduces the information rate, that is, the density of the public knapsack,
which makes lattice-based attacks more relevant rather than less. Iteration
therefore did not repair the Merkle--Hellman idea. It only changed the form of
the weakness and confirmed that simple repeated transformations were not enough
to restore security.~\cite[p.~344]{Brickell1984}\cite[p.~8]{OdlyzkoRiseFall}

\subsection{Chor--Rivest Knapsack-Type Cryptosystem}

The Chor--Rivest cryptosystem was a more substantial attempt to build a
knapsack-type public-key system after the weaknesses of Merkle--Hellman
variants became clear. It does not hide a superincreasing sequence by modular
multiplication. Instead, it uses arithmetic in finite fields and a construction
related to the Bose--Chowla theorem. One of the aims of this design was to
obtain a knapsack with higher density, because low-density lattice attacks are
most effective when the public weights are large compared with the number of
message bits.~\cite[pp.~901--902]{ChorRivest1988}

The construction works over a finite field $GF(p^h)$. Let $t$ be an element of
degree $h$ over $GF(p)$, and let $g$ be a generator of the multiplicative group
of $GF(p^h)$. For each $\alpha \in GF(p)$, a value $a_\alpha$ is defined by
\[
    g^{a_\alpha}=t+\alpha.
\]
The public weights are obtained by applying a secret permutation $\pi$ and
adding a secret offset $d$:
\[
    c_i \equiv a_{\pi(i)} + d \pmod{p^h-1}.
\]
The public key contains the values $c_i$ together with the public parameters,
while the trapdoor information contains the finite-field structure, the
generator, the permutation, and the offset.~\cite[pp.~902--903]{ChorRivest1988}

Encryption still has the form of a knapsack sum. A plaintext block is encoded
as a binary vector $m=(m_0,\dots,m_{p-1})$ with exactly $h$ ones, so
\[
    \sum_{i=0}^{p-1} m_i = h.
\]
The ciphertext is computed as
\[
    S \equiv \sum_{i=0}^{p-1} m_i c_i \pmod{p^h-1}.
\]
From the public point of view, encryption therefore still resembles a subset-sum
problem. The difference is in the trapdoor used for decryption.

To decrypt, the legitimate receiver first removes the offset by computing
$S-hd$ modulo $p^h-1$. Exponentiating $g$ to this value gives a product of
terms of the form $t+\alpha$. This product can be represented as a polynomial
over $GF(p)$, and the receiver recovers the message by factoring it into linear
terms. This is different from Merkle--Hellman decryption, where the receiver
returns to a superincreasing sequence and applies the greedy algorithm.
Chor--Rivest is therefore better understood as a different knapsack-type idea,
not as a simple modification of the original Merkle--Hellman construction.~\cite[pp.~903--904]{ChorRivest1988}

At the time of Odlyzko's survey, Chor--Rivest was described as one of the few
knapsack-type systems that had resisted the attacks known then.~\cite{OdlyzkoRiseFall}
This statement has to be understood historically. Later cryptanalysis by
Vaudenay showed how to break the proposed Chor--Rivest parameters, including
parameters based on $GF(p^{24})$ and $GF(256^{25})$.~\cite[pp.~243,
255]{Vaudenay1998} For this thesis, the importance of Chor--Rivest is mainly
that it shows how researchers tried to move away from low-density
superincreasing constructions toward high-density knapsack-type designs.

\subsection{Other Historical Proposals (Brief Overview)}

Several other knapsack-based constructions were proposed during the early
development of public-key cryptography. Merkle and Hellman also described a
multiplicative trapdoor knapsack, where the trapdoor is based on products of
relatively prime integers rather than on a superincreasing additive
sequence.~\cite[pp.~526--527]{MerkleHellman1978} In this construction, the
public values are related to discrete logarithms, and decryption uses the
secret parameters to transform the ciphertext into a product that can be
factored easily by the legitimate receiver. This idea tried to combine the
knapsack problem with number-theoretic structure, but Odlyzko later showed that
simple variants of the multiplicative knapsack are also vulnerable under
practical parameter choices.~\cite{Odlyzko1984Attacks}

Knapsack ideas were also proposed for authentication and digital signatures.
One important example was Shamir's fast signature scheme, which was designed to
provide efficient signing rather than encryption. Odlyzko showed that this
scheme could be attacked quickly using lattice-based methods, without
recovering the exact original secret information. He also noted that similar
methods could be adapted to the Schobi--Massey knapsack-based authentication
scheme. These examples show that the weaknesses of knapsack constructions were
not limited to message encryption, but also affected related signature and
authentication proposals.~\cite{Odlyzko1984Attacks}

Other proposals, such as Graham--Shamir knapsacks, generalized knapsacks, and
so-called ultimate knapsack variants, also appeared as attempts to repair or
extend the original idea. They are not described in detail here, because their
constructions differ in technical details and would require separate analysis.
For the purpose of this thesis, the important point is that many of these
systems were attacked by similar families of methods: structural
reconstruction, Diophantine approximation, and lattice basis reduction. This
broader pattern supports the main lesson of the chapter: changing the surface
form of a knapsack cryptosystem did not necessarily remove the hidden structure
that made cryptanalysis possible.~\cite{Brickell1984,OdlyzkoRiseFall}

\subsection{Efficiency and Implementation Complexity}
\subsection{Summary Table of Schemes and Attacks}
